% ============================================================
%  第9章 案例分析：图像聚类
%  源文件：content/chapter9/
% ============================================================

\emph{词袋}(BoW，bag-of-words)模型是大多流行图像分类模型中的一个，这个模型是图像搜索系统的一个组件。BoW模型将图片的特征作为一组词，用这组词作为一个矢量用来描述这幅图像。本章我们会讨论如何使用OpenCL实现BoW模型中重要的组件——\emph{建立直方图}(Histogram Builder)，以及学习如何对OpenCL源代码进行优化。


% ---- 9.1-chinese.md ----
\section{图像聚类简介}

图像分类源于机器视觉，其根据图像可见内容对图像进行分类。例如，某个图像算法可能会用来告诉你该图片中是否有人。虽然检测人可能是很简单的事情，但能将图像进行准确分类的算法，仍然是目前的所要面临的挑战。

BoW模型通常会用于文本分类或自然语言处理。BoW模型中每个词出现的频度都会作为一个训练参数传入对应的机器学习算法中。除了进行文本分类，BoW也可以应用于图像。为了让BoW能对图像进行分类，我们需要从图像中获取一组词(和文本分类一样)，以及对应词的出现频率。机器视觉中，这里的“词”通常被看做是图像的一个“特征”。特征算法将图像缩小成一组特征值。图9.1展示了个图像分类的高端算法，其包括生成特征、聚合和直方图建立这几步。

\myGraphic{0.92}{images/chapter9/9-1.png}{图9.1 图像分类流程。像SURF这样的算法可以用来产生特征。像k-means类的聚合算法之后会产生一组质心特征，这些特征可以用来描述原始图像。通过直方图构建算法将这些计算出的特征赋予质心。}

\myGraphic{0.92}{images/chapter9/9-2.png}{图9.2 使用SURF算法生成特征。SURF算法可以将一副图像作为输入，并且产生一组对应的特征。每个特征包括图像位置信息和一个具有64个值的描述符。}

下面简单的描述一下这几步：

\begin{enumerate}
\item 生成特征：我们已经将特征产生算法应用于BoW模型，将其称为\emph{加速稳健特征算法}(SURF，Speeded Up Robust Features)。SURF算法在2006年，由Bay等人在欧洲计算机视觉会议(ECCV)上发表[1]，其对不同大小的图像进行尺度不变的变换。如图9.2所示，将图像出入SURF算法中，将会得到一组用来描述原始图像的特征值。每个特征值包含图像位置和描述矢量。每个特征中的描述符向量是一个64维的向量。没描述符中包括该特征颜色和特征位置周围颜色梯度关系。本章中的特征，都具有一个有着64个元素的描述符。其余的成员变量则不是本章关注的焦点。
\item 图像聚合：SURF产生的描述符通常都经过量化，通常会经过k-means聚合，然后映射到群组中。每个群组的质心就是所谓的“视觉词”。
\item 直方图构建：这一阶段的任务是将SURF算法产生的描述符设置到直方图的视觉词中去。为了完成这项工作，我们需要确定描述符和质心的对应关系。本例中，SURF产生的特征描述符和质心都由64维构成。我们通过计算描述符和所有质心的欧式距离，并且将描述符赋予距离最近的质心。直方图在这里就是用于统计每个质心被赋予描述符的次数。
\end{enumerate}

使用这种方法，我们就能表示视觉词出现的频度。机器算法，例如支持向量机就能使用模型对图片进行分类。这一过程在图9.1中展示。

本章我们将探索在直方图构建步骤上的并行化。我们首先使用在CPU上的串行算法对该算法做一个具体的了解，然后使用OpenMP进行CPU版本的并行化优化。然后将多线程实现换成OpenCL，让这一阶段的计算在GPU上完成。我们会有多个GPU版本的实现，包括原始实现，对全局内存进行优化的实现，以及利用局部内存的实现。本章最后，我们将使用AMD的Radeon HD 7970 GPU对相应的实现进行性能评估。

\begin{myNotic}{贴士}
本章着重于图像分类算法中的直方图构建步骤。读者如果对特征生成算法感兴趣，这里提供一个开源版本的SURF实现(OpenCL)：\texthref{https://code.google.com/p/clsurf/} google code无法使用，GitHub地址：\texthref{https://github.com/perhaad/clsurf}
\end{myNotic}

\myRefRule

\begin{myRefs}

[1] H. Bay, T. Tuytelaars, L.V.Gool, SURF:speed up robust features, in: ECCV, 2006, pp.404-417.

\end{myRefs}



% ---- 9.2-chinese.md ----
\section{直方图的特性——CPU实现}

本节我们将介绍如何将SURF计算出的特征转换成直方图，我们先用CPU是现实一个串行执行的版本。然后使用OpenMP使用CPU多核来完成这个算法的并行化。

\subsection{串行实现}

\begin{codeboxt}[cpp]{代码清单9.1 将SURF特征设置到集群直方图的串行版本}
// Loop over all the descriptors generated for the image
for (int i = 0; i < n_des; i++){
  membership = 0;
  min_dist = FLT_MAX;
  // Loop over all cluster centroids available
  for (j = 0; j < n_cluster; j++){
    dist = 0;
    // n_featrues: No. of elements in each descriptor (64)
    // Calculate the distance between the descriptor and the centroid
    for (k = 0; k < n_features; k++){
      dist_temp = surf[i][k] - cluster[j][k];
      dist += dist_temp * dist_temp;
    }
    // Update the minimum distance
    if (dist < min_dist){
      min_dist = dist;
      membership = j;
    }
  }
  // Update the histogram location of the closest centroid
  histogram[membership] += 1;
}
\end{codeboxt}

代码清单9.1中，展示了如何将SURF计算出来的特征设置为集群的质心(视觉词)。第2行遍历了每个SURF特征的描述符，第7行遍历了集群的所有质心。第12行循环遍历当前描述符中的64个元素，并计算当前特征与当前集群质心的欧氏距离。第18行找到离集群质心最近的SURF特征，并将其设置为成员。

\subsection{OpenMP并行实现}

为了展现CPU多核多线程的能力，我们使用OpenMP来对清单9.1的代码进行并行化。OpenMP的编程接口支持多平台的内存共享并行编程，可以使用C/C++和Fortran作为编程语言。其定义了一种可移植、可扩展的简单模型，并且灵活的接口可以让当前的应用立即化身为多线程应用[2]。在C/C++代码中，OpenMP使用编译标识(\#pragma)直接告诉编译器生成对应的多线程实现。

\begin{codeboxt}[cpp]{代码清单9.2 将清单9.1的代码进行多核并行化}
#pragma omp parallel for schedule(dynamic)
// Loop over all the descriptors generated for the image
for (int i = 0; i < n_des; i++){
  membership = 0;
  min_dist = FLT_MAX;
  // Loop over all cluster centroids available
  for (j = 0; j < n_cluster; j++){
    dist = 0;
    // n_featrues: No. of elements in each descriptor (64)
    // Calculate the distance between the descriptor and the centroid
    for (k = 0; k < n_features; k++){
      dist_temp = surf[i][k] - cluster[j][k];
      dist += dist_temp * dist_temp;
    }
    // Update the minimum distance
    if (dist < min_dist){
      min_dist = dist;
      membership = j;
    }
  }
  // Update the histogram location of the closest centroid
  #prargma omp atomic
  histogram[membership] += 1;
}
\end{codeboxt}

使用OpenMP可以将直方图构建任务分配到多个CPU核上。每个线程处理不同的描述符和集群执行的距离，并将相应的描述符赋予质心。虽然每个描述符的计算是独立的，但是将值赋予质心的过程会出现条件竞争：多个线程想要同时更新同一个位置的内存，那么结果将是不可预测的。条件竞争可以使用原子加操作进行解决(第26行)。

清单9.2中，编译标识出现在第2行，其作用就是将每一次循环迭代放置到一个线程中。第18行的标识则告诉编译器，使用原子操作来更新共享内存。

\myRefRule

\begin{myRefs}

[2] L. Dagum, R. Menon, OpenMP: an industry standard API for shared-mempry programming, IEEE Comput. Sci. Eng. 5(1)(1998)46-55

\end{myRefs}



% ---- 9.3-chinese.md ----
\section{OpenCL实现}

本节，我们将了解一下如何在GPU上实现直方图构建。首先实现实现了一个常规内核，该内核代码基于串行和OpenMP版本的算法。然后，再来对GPU实现进行优化，这里使用的优化策略为：合并访问和局部内存。

\subsection{常规GPU实现：GPU1}

根据清单9.2中的OpenMP实现，我们可以实现一个最简单OpenCL版本，也就是将循环迭代拆开，让每个工作项完成一个标识符的归属计算。

不过，OpenMP实现中，我们遇到了条件竞争的问题，这里工作项需要对在全局内存上的直方图进行更新。为了解决这一问题，在OpenCL中我们依旧使用原子操作。下面的OpenCL内核就是我们的实现。我们暂且称这段代码为GPU1：

\begin{codeboxt}[cpp]{代码清单9.3 内核GPU1，直方图基线内核}
__kernel
void kernelGPU1(
  __global float *descriptors,
  __global float *centroids,
  __global int *histogram,
  int n_descriptors,
  int n_centroids,
  int nfeatures){

  // Global ID identifies SURF descriptor
  int desc_id = get_global_id(0);

  int membership = 0;
  float min_dist = FLT_MAX;

  // For each cluster, compute the membership
  for (int j = 0; j < n_centroids; j++){
    float dist = 0;

    // n_features: No. of elements in each descriptor(64)
    // Calculate the distance between the descriptor and the centroid
    for (int k = 0; k < n_features; k++){
      float temp = descriptors[desc_id * n_features + k] -
        centroids[j * n_features + k];
      dist += temp * temp;
    }

    // Update the minimum distance
    if (dist < min_dist){
      min_dist = dist;
      membership = j;
    }
  }

  // Atomic increment of histogram bin
  atomic_fetch_add_explicit(&histogram[membership], 1, memory_order_relaxed, memory_scope_device);
}
\end{codeboxt}

注意清单9.3中，第37行的原子加操作，使用的自由序进行。我们选择这个内存序的原因是，更新操作比较简单，并不需要对内存访问顺序进行严格要求。详细内容可回顾第7章中的相关章节。

\subsection{合并访问内存：GPU2}

例子中对于数据的访问不存在跨距，所以工作项可以在GPU执行单指令多数据命令(SIMD)。并且SURF描述符和集群质心矢量都具有64个连续数据。再来看看清单9.3中的第24行，并且注意描述符访问的方式。对于给定的GPU硬件来说，第24行对于内存的访问是否高效呢？

假设有4个并行的工作项，其全局索引分别是0到3。在GPU执行最内部的循环时，这个四个工作项所访问的数据间具有很大的跨距——在这个版本的内核代码中，数据跨距为n\_\allowbreak{}features。假设我们现在在处理质心0(j=0)，并且正在处理质心0与特征0(k=0)之间的距离。那么工作项分别访问的数据为：descriptors[0]、descriptors[64]、descriptors[128]、descriptors[192]。那么计算下一个特征时就要访问descriptors[1]、descriptors[65]、descriptors[129]、descriptors[193]，以此类推。

第8章中我们介绍了对于连续数据的合并访问方式，该方式只向内存系统发出更少的请求，以高效的方式获取相应的数据。不过，跨距访问方式需要产生多个访存请求，从而导致性能下降。这种带有跨距的访存方式是无法进行合并访问的。

为了通过使用合并访问的方式提升内存带宽的利用率，我们需要调整数据在descriptors存放的顺序。可以通过一种常用的矩阵操作来完成——转置，如图9.3所示。转置过程中，矩阵的行列坐标互换。我们可以创建一个简单的内核来完成这项工作。

\myGraphic{0.92}{images/chapter9/9-3.png}{图9.3 将描述符中的数据进行转置，以便进行合并访问}

\myGraphic{0.92}{images/chapter9/9-4.png}{图9.4 对一维数组进行转置。}

一维数组的转置如图9.4所示。转置后，descriptors[0]、descriptors[64]、descriptors[128]、descriptors[192]就是连续的。并且内核内部可以进行合并访问，只需要一个访存请求，即可直接将4个连续的数据取出。

\begin{codeboxt}[cpp]{代码清单9.4 使用内存合并访问的内核——GPU2}
__kernel
void kernelGPU2(
  __global float *descriptors,
  __global float *centroids,
  __global int *histogram,
  int n_descriptors,
  int n_centroids,
  int nfeatures){

  // Global ID identifies SURF descriptor
  int desc_id = get_global_id(0);

  int membership = 0;
  float min_dist = FLT_MAX;

  // For each cluster, compute the membership
  for (int j = 0; j < n_centroids; j++){
    float dist = 0;

    // n_features: No. of elements in each descriptor(64)
    // Calculate the distance between the descriptor and the centroid
    for (int k = 0; k < n_features; k++){
      float temp = descriptors[k * n_descriptors + desc_id] -
        centroids[j * n_features + k];
      dist += temp * temp;
    }

    // Update the minimum distance
    if (dist < min_dist){
      min_dist = dist;
      membership = j;
    }
  }

  // Atomic increment of histogram bin
  atomic_fetch_add_explicit(&histogram[membership], 1, memory_order_relaxed, memory_scope_device);
}
\end{codeboxt}

清单9.4中的第24行，现在使用的是k * n\_\allowbreak{}descriptors + desc\_\allowbreak{}id，可与清单9.3进行对比。当k和n\_\allowbreak{}descriptors具有相同的值时，所有工作项就挨个对数据进行计算，不同工作项可以通过唯一标识符对数据进行访问(desc\_\allowbreak{}id)。我们之前举了4个工作组的例子，在k=0时，就只需要访问descriptors[0]、descriptors[1]、descriptors[2]和descriptors[3]即可。当k=1时，则需要访问descriptors[64]、descriptors[65]、descriptors[66]和descriptors[67]。这样的访存方式是最理想的，并且在GPU上进行合并访问可以高效的对内存系统进行操作。

\subsection{向量化计算：GPU3}

当SURF算法生成特征之后，描述符的长度固定为64维，这时使用向量float4可以对应用进行提速。在CPU上，对于其扩展流媒体SIMD指令来说，能够获得较高的吞吐量。同样的方式也适用于GPU(比如，AMD Radeon 6xxx 系列)，在GPU上使用向量化方式的计算也能带来相应的性能收益。AMD和NVIDIA的新型GPU都不会显式的执行向量指令，不过在某些情况下向量的方式也会提升内存系统的使用效率。

\begin{codebox}[cpp]
float a[4], b[4], c[4];

a[0] = b[0] + c[0];
a[1] = b[1] + c[1];
a[2] = b[2] + c[2];
a[3] = b[3] + c[3];
\end{codebox}

向量化可以推广到一些可扩展的数据类型，比如：向量、矩阵和高维数组。这里的例子中，我们将数组中的元素由编程者显式的放入float4类型中。下面的代码就向量化所要进行的一些操作。

\begin{codebox}[cpp]
float a[4], b[4], c[4];

float4 b4 = (float4)(b[0], b[1], b[2], b[3]);
float4 c4 = (float4)(b[0], b[1], b[2], b[3]);
float4 a4 = b4 + c4;
\end{codebox}

为了将向量化应用到我们的算法中，我们在这里更新一下内核实现(代码清单9.5)，称这个内核为GPU3。

\begin{codeboxt}[cpp]{代码清单9.5 使用向量化的内核代码——GPU3}
__kernel
void kernelGPU3(
  __global float *descriptors,
  __global float *centroids,
  __global int *histogram,
  int n_descriptors,
  int n_centroids,
  int nfeatures){

  // Global ID identifies SURF descriptor
  int desc_id = get_global_id(0);

  int membership = 0;
  float min_dist = FLT_MAX;

  // For each cluster, compute the membership
  for (int j = 0; j < n_centroids; j++){
    float dist = 0;

    // n_features: No. of elements in each descriptor(64)
    // Calculate the distance between the descriptor and the centroid
    // The increment of 4 is due to the explicit verctorization where
    // the distance between 4 elements is calculated in each
    // loop iteration
    for (int k = 0; k < n_features; k++){
      float4 surf_temp = (float4)(
        descriptors[(k + 0) * n_descriptors + desc_id],
        descriptors[(k + 1) * n_descriptors + desc_id],
        descriptors[(k + 2) * n_descriptors + desc_id],
        descriptors[(k + 3) * n_descriptors + desc_id]);
      float4 cluster_temp = (float4)(
        centroids[j * n_feature + k],
        centroids[j * n_feature + k + 1]
        centroids[j * n_feature + k + 2]
        centroids[j * n_feature + k + 3]);

      float4 temp = surf_temp - cluster_temp;
      temp = temp * temp;

      dist += temp.x + temp.y + temp.z + temp.w;
    }

    // Update the minimum distance
    if (dist < min_dist){
      min_dist = dist;
      membership = j;
    }
  }

  // Atomic increment of histogram bin
  atomic_fetch_add_explicit(&histogram[membership], 1, memory_order_relaxed, memory_scope_device);
}
\end{codeboxt}

\subsection{将SURF特征放入局部内存：GPU4}

下面是代码清单9.4中，访问descriptors和centroids的代码片段：

\begin{codebox}[cpp]
for (int k = 0; k < n_features; k++){
  float temp = descriptors[k * n_features + k] -
    centroids[j * n_features + k];
  dist += temp * temp;
}
\end{codebox}

从上面的片段可以看出，要对这个内存进行多次的访问。有没有可能将这段内存放置到OpenCL特定的内存空间中，以提升性能呢？

这里访问centroids的位置可以使用工作项的索引进行表示。常量内存非常适合以这种方式进行访问，下节中我们将介绍如何将centroids放入常量内存。当前版本中，重点放在对descriptors访存的优化。当使用工作项索引对descriptors数据进行访存时，因为其存在访存跨度的原因，所以不是很适合放入常量内存中。那么这个数组能放入到局部内存中吗？

GPU上的局部内存是一段具有高带宽、低延迟的内存区域，其可以将数据共享给工作组内每一个工作项。GPU上有专用的局部内存，访问局部内存的速度通常要比全局内存快很多。同样，与全局内存访问不同，访问局部内存通常都不需要合并访问，就算是在局部内存上产生了内存访问冲突，其性能也要优于全局内存。不过，局部内存的大小有限——在AMD Radeon HD 7970 GPU上每个计算单元只有64KB大小的局部内存，所以能分配给一个每个工作组的只有32KB。若是为每个工作组分配一个很大的内存，则会限制GPU上执行线程的数量。对于GPU来说，减少线程就意味着不能很好的掩盖访存延迟，同样也会让计算资源闲置。

最初，数据放到局部内存上似乎并不是最好的选择，因为局部内存能在工作项间共享数据，不过对于descriptors来说没有数据是需要共享的。不过，对于很多GPU，局部内存具有一些额外的优势。首先，局部内存与本地数据存(LDS)相对应，其比L1缓存大4倍。其次再说缓存命中，LDS内存访存的延迟也要比L1缓存低很多。为了充分复用，数据放在LDS将会带来比放在L1缓存上更好的性能，因为LDS的延迟很小，即使是高频度访问也能轻松应对。这里所要权衡的是局部内存使用的多少与执行工作组数量，换句话说就是更多的使用执行单元还是内存系统。优化策略的权衡，需要根据不同的架构特点决定。代码清单9.6展示了，descriptors使用局部内存的内核代码。

\begin{codeboxt}[cpp]{代码清单9.6 将descriptor数据存放在局部内存中的内核——GPU4}
__kernel
void kernelGPU4(
  __global float *descriptors,
  __global float *centroids,
  __global int *histogram,
  int n_descriptors,
  int n_centroids,
  int nfeatures){

  // Global ID identifies SURF descriptor
  int desc_id = get_global_id(0);
  int local_id = get_local_id(0);
  int local_size = get_local_size(0);

  // Store the descriptors in local memory
  __local float desc_local[4096]; // 64 descriptors * 64 work-items
  for (int i = 0; i < n_features; i++){
    desc_local[i * local_size + local_id] =
        descriptors[i * n_descriptors + desc_id];
  }
  barrier(CLK_LOCAL_MEM_FENCE);

  int membership = 0;
  float min_dist = FLT_MAX;

  // For each cluster, compute the membership
  for (int j = 0; j < n_centroids; j++){
    float dist = 0;

    // n_features: No. of elements in each descriptor(64)
    // Calculate the distance between the descriptor and the centroid
    for (int k = 0; k < n_features; k++){
      float temp = descriptors[k * local_size + desc_id] -
        centroids[j * n_features + k];
      dist += temp * temp;
    }

    // Update the minimum distance
    if (dist < min_dist){
      min_dist = dist;
      membership = j;
    }
  }

  // Atomic increment of histogram bin
  atomic_fetch_add_explicit(&histogram[membership], 1, memory_order_relaxed, memory_scope_device);
}
\end{codeboxt}

访问descriptor时，n\_\allowbreak{}descriptors和desc\_\allowbreak{}id对于所有工作项来说都是固定的。索引值都是基于k变化的。因此，每个工作项都能访问到desctiptor中的n\_\allowbreak{}feature个元素(64)。GPU上L1缓存的大小非常小，每次缓存访问丢失都会给访问全局内存的过程带来更多冗余的开销。

将secriptors放入LDS将需要64x4=256字节。一个波面阵中有64个工作项，每个波面阵就需要是用16KB的LDS空间用来缓存desctiptor。当有64KB的LDS空间时，就能在每个计算单元上运行4个波面阵(每个波面阵只有一个工作项)。在HD 7970上，每个计算单元由4个SIMD单元组成，这样的话每个SIMD单元就只能处理一个波面阵，SIMD单元之间的延迟掩盖就没有了。为了获取最佳性能，我们需要在低延迟访问和减少并行化中进行权衡。

\subsection{将聚类中点坐标放入常量内存：GPU5}

第4章中我们讨论过卷积操作，并且在第7章介绍过内存模型，常量内存所放置的数据可以让所有工作项同时访问。通常放置在常量内存中的数据有卷积的内核和一些常量数据。直方图的例子中，每个质心的描述符就是固定的数据。访问centroids时，地址要根据两层循环的索引进行计算得出，并非按照工作项索引。因此，所有工作项都要进行一定的计算才能产生所要访问的地址。

GPU常量内存通常会映射到一块较为特殊的缓存硬件上，其大小是固定的。所以，将centroids放置到常量内存中时，需要考虑其数据的大小。本例中Radeon HD 7970，其常量内存大小为64KB。本例中每个特征的质心数据有256字节。因此，我们最多只能同时将256个质心放置到常量内存中。

相信会有读者对映射到常量内存上数据的性能有疑问。如果所有工作项同时访问同一个地址，以非合并访问的方式下，是否对常量内存只产生一个访问请求？大多数现代GPU上，对这个问题的回答是“是”：只产生一个访问请求。不过，与LDS上的descriptors类似，将centroids放置在常量内存上也会带来性能收益。首先，使用常量内存会减小GPU上L1缓存的压力。GPU上的L1缓存比较小，现在将一块64KB的空间移除L1缓存，将其他数据放置在这里，从而会带来巨大的性能提升。其次，GPU常量内存的访存延迟也要比通用L1缓存低很多。当我们的数据不是那么多的时候，我们可以考虑将其放置在常量内存中，这将带来很好的性能收益。

将centroids数组映射到常量内存的方式很简单，只需要将\texttt{\_\_global}改成\texttt{\_\_constant}即可。具体代码见代码清单9.7。

\begin{codeboxt}[cpp]{代码清单9.7 将质心数据放置在常量内存中——GPU5}
__kernel
void kernelGPU4(
  __global float *descriptors,
  __constant float *centroids,
  __global int *histogram,
  int n_descriptors,
  int n_centroids,
  int nfeatures){

  // Global ID identifies SURF descriptor
  int desc_id = get_global_id(0);
  int local_id = get_local_id(0);
  int local_size = get_local_size(0);

  // Store the descriptors in local memory
  __local float desc_local[4096]; // 64 descriptors * 64 work-items
  for (int i = 0; i < n_features; i++){
    desc_local[i * local_size + local_id] =
        descriptors[i * n_descriptors + desc_id];
  }
  barrier(CLK_LOCAL_MEM_FENCE);

  int membership = 0;
  float min_dist = FLT_MAX;

  // For each cluster, compute the membership
  for (int j = 0; j < n_centroids; j++){
    float dist = 0;

    // n_features: No. of elements in each descriptor(64)
    // Calculate the distance between the descriptor and the centroid
    for (int k = 0; k < n_features; k++){
      float temp = descriptors[k * local_size + desc_id] -
        centroids[j * n_features + k];
      dist += temp * temp;
    }

    // Update the minimum distance
    if (dist < min_dist){
      min_dist = dist;
      membership = j;
    }
  }

  // Atomic increment of histogram bin
  atomic_fetch_add_explicit(&histogram[membership], 1, memory_order_relaxed, memory_scope_device);
}
\end{codeboxt}



% ---- 9.4-chinese.md ----
\section{性能分析}

为了展示不同的内核实现对于性能的影响，我们将这些内核都在Radeon HD 7970 GPU上执行了一遍。为了展示对不同数据大小的性能优化，我们也生成除了对应的SURF描述符和集群质心。生成的SURF描述符数量为4096，16384和65536。同时，对应质心的数量为16,64和256。我们选取数量较大的SURF特征，是因为对于一张高分辨率的图来说，通常都包含成千上万个特征。不过，对于质心来说其数量就相对较少，如果质心的数量过多会影响图像分类的准确度。

这里的性能数据是由相应的性能分析工具给出，这里我们使用了AMD公司的CodeXL，第10章中我们会对该工具进行介绍。本章主要描述对于OpenCL内核代码进行优化后的性能情况。这里还需要注意的是，目标平台的架构对于优化代码来说也有很大的影响。

\subsection{GPU性能}

我们将GPU1内核作为OpenCL内核实现的基线。在使用GPU2处理描述符和质心前，我们是用到了矩阵转置对描述符进行了转置操作。转置操作属于另外一个内核，为了避免这个因素影响我们之后优化的内核，这里将转置的时间单独列出来，在表9.1中展示。那么第二个内核，也就是直方图内核的性能数据在表9.2中展示。

表9.1 转置内核执行的时间

\par\medskip\noindent
\begingroup\fontsize{9bp}{12.5bp}\selectfont\renewcommand{\arraystretch}{1.12}
\begin{center}
\begin{tabular}{ll}
\toprule
特征数量 & 转置内核耗时(ms) \\
\midrule
4096 & 0.05 \\
16,384 & 0.50 \\
65,536 & 2.14 \\
\bottomrule
\end{tabular}
\end{center}
\endgroup
\medskip

表9.2 直方图——不同内核实现的耗时统计

\par\medskip\noindent
\begingroup\fontsize{9bp}{12.5bp}\selectfont\renewcommand{\arraystretch}{1.12}
\begin{center}
\begin{tabular}{lllllll}
\toprule
集群数量 & SURF描述符数量 & GPU1 & GPU2 & GPU3 & GPU4 & GPU5 \\
\midrule
8 & 4096 & 0.41 & 0.27 & 0.10 & 0.17 & 0.09 \\
8 & 16,348 & 3.60 & 0.28 & 0.17 & 0.69 & 0.19 \\
8 & 65,536 & 15.36 & 1.05 & 0.59 & 1.31 & 0.74 \\
16 & 4096 & 0.77 & 0.53 & 0.19 & 0.28 & 0.14 \\
16 & 16,348 & 7.10 & 0.53 & 0.32 & 0.57 & 0.29 \\
16 & 65,536 & 30.41 & 1.47 & 1.17 & 2.26 & 1.12 \\
64 & 4096 & 6.00 & 3.53 & 1.34 & 1.00 & 0.43 \\
64 & 16,348 & 28.28 & 2.11 & 1.20 & 2.96 & 0.86 \\
64 & 65,536 & 122.09 & 5.80 & 4.65 & 9.04 & 3.87 \\
128 & 4096 & 4.96 & 4.04 & 1.47 & 1.95 & 0.81 \\
128 & 16,348 & 55.70 & 4.27 & 2.40 & 5.89 & 1.61 \\
128 & 65.536 & 243.30 & 11.63 & 9.29 & 17.46 & 6.43 \\
256 & 4096 & 10.49 & 8.06 & 2.84 & 4.35 & 1.57 \\
256 & 16,348 & 109.67 & 8.62 & 4.77 & 11.44 & 3.13 \\
256 & 65,536 & 488.54 & 23.28 & 18.71 & 34.73 & 13.97 \\
\bottomrule
\end{tabular}
\end{center}
\endgroup
\medskip



% ---- 9.5-chinese.md ----
\section{本章总结}

本章在实际环境中执行了OpenCL的多种优化内核。优化包括数据转化，向量化的数学操作，以及将数据映射到局部和常量内存上。我们也通过数据观察到，不同的优化措施对于性能的影响。


